\chapter{Is Taste Really Personal?}
\section{The Playlist in Your Headphones}
First, pick up something you probably have with you: a pair of headphones.

You need not actually take them out; doing it in your head will do. Imagine putting them on, opening your music app, and selecting the playlist you have heard hundreds of times. The first song's introduction begins, and your shoulders relax without your noticing. This playlist is yours---the app knows it, you know it, and friends visiting your profile can see it. Like your fingerprints, it seems unique: surely nobody else in the world likes precisely these thirty-odd songs in precisely this order.

Now do something a little uncomfortable. Scroll to the top, look through the playlist, and ask yourself: how did each of these songs get here in the first place?

The first may come from the CD Dad always played while driving when you were in primary school. You barely noticed it then, but ten years later, late one night, you suddenly remembered it, tracked it down, and added it. Listening brought a lump to your throat. The second was the song everyone in seventh grade listened to; without knowing it, you could not join the conversation at break. The third came from a short video: you heard only fifteen seconds of its climax and do not even know how the introduction goes. The fourth was shared by an older student you admired. It did little for you, but you saved it anyway. The fifth may genuinely be your own discovery---something shuffle played late one night, when you were alone in your room and it made your scalp tingle.

Look: every song on this supposedly most personal playlist has a history. Some came from family, some from peers, some from a server you have never seen, and some from the person you want to become. Very few have no discernible origins, belonging entirely to an encounter between you and the music.

That is this chapter's question: is taste really personal? When you say, ``I like it,'' how many other people inhabit the ``I'' doing the talking?

\section{Judgment in a Ten-Minute Break}
Now come with me into a scene.

Time: a morning in May, the ten-minute break after the second lesson. Place: an eighth-grade classroom. A ceiling fan turns slowly overhead with a steady hum. Chalk writing remains partly unerased at the front, a pile of collected exercise books sits on the windowsill, and outside comes the whistle from another class's PE lesson.

A girl in the back row, Xiao Ai, puts her phone on the corner of her desk and plays a song aloud. It is a recent short-video hit. The introduction begins with bright, busy drumbeats, and the chorus has an irresistibly catchy vocal flourish. She hums along, tapping her foot under the desk.

A Zhe, sitting diagonally ahead, turns around. He is the kind of boy who quotes lyrics in his essays; his headphones always play bands you have never heard of, with black-and-white covers and long names. He listens for a few seconds, then laughs, just loudly enough for three people nearby to hear:

``You listen to this too?''

Three characters and a question mark. Someone nearby laughs along. Xiao Ai's hand freezes in midair, and her face slowly reddens. She wants to say something---``What's it to you?'' perhaps---but swallows the words, because A Zhe immediately twists the knife: ``These songs come off an assembly line. They can't even be bothered to change the chords. Hear two and you've heard them all.''

Xiao Yu, beside her, had been about to ask the song's name. Hearing this, she quietly pulls her own phone into her sleeve: more than half her playlist consists of songs like this.

The ten-minute break ends, the bell rings, and outwardly the incident is over. But something remains: Xiao Ai plays no more songs aloud that day; Xiao Yu makes her playlist private after getting home; and A Zhe shares his favorite band's new album that evening with a three-word caption: ``Cleanse your ears.''

Remember this classroom. Everything else in this chapter addresses the question hidden in those ten minutes: what gives A Zhe the right? And does Xiao Ai really have anything to be ashamed of?

\section{Is ``I Just Like It'' Enough?}
Let us lay out the conflict without rushing to judgment.

Xiao Ai has a simple but sturdy argument. Liking is liking: it grows in me, and any hurt is mine to feel. This song makes me happy, and that happiness is real. I walk more lightly when I hum it. Is that not evidence? You need no qualification to like a person, so why need one to like a song? I do not understand A Zhe's talk of chords and assembly lines, but I need not understand it. Must happiness first pass a music-theory examination to count?

A Zhe's argument sounds strong too. There really is a difference between something made with care and something dashed off. A song that gets by on three chords and one whose arrangement contains a dozen layers of detail sit on the same shelf. A careless buyer may not distinguish them, but failing to distinguish them does not erase the difference. He says he is not targeting Xiao Ai; he is defending music. If everyone listens only to fifteen-second choruses, serious musicians will starve, until fifteen-second choruses are all the shelf contains.

Behind these positions lies a question far deeper than who lacks social skills:

\textbf{Are your preferences actually your own?}

Notice that this question points both ways. Toward A Zhe: you call your taste superior, but where did it come from? Have you heard those bands because you have time, equipment, and a cousin who introduced you, or because you were born superior? Had you grown up in Xiao Ai's family, can you guarantee you would listen to the same things? Toward Xiao Ai: you say, ``I just like it,'' but how did the song reach your ears? Did you freely select it from millions, or did a platform calculate that you would linger and push it into your fifteen seconds? Your happiness is real, but you did not build the road that brought you to the song.

Worse, the question will not settle on either side. If all taste is shaped, A Zhe's superiority is shaped too, leaving him nothing to boast about. If taste is purely personal, are platforms and businesses spending so much to study your preferences out of charity?

Before reading further, take a position in that classroom. Does A Zhe's ``You listen to this too?'' contain even a grain of truth?

\section{Where Preferences Come From: Five Hands}
To ask whether a preference is your own, first examine how it was installed in your head. Let us inspect the five hands reaching into your playlist. The point is not to declare you manipulated, but to make the process visible.

\textbf{The first hand: family.}

A person's earliest musical memories are almost never chosen: Dad's old driving songs, Mom's cooking-time humming, opera on Grandma's radio. In many Japanese families, children grow up immersed in elders' enka, a traditional popular-song style with slow melodies and pronounced vibrato. Many Mexican children grow up with mariachi played by street bands at birthday parties. Before you can choose, these sounds establish a baseline for what sounds pleasant and what sounds noisy. Psychologists have identified a simple pattern: repeated exposure can make people like something. Familiarity itself is a sugar coating. Much of what seems an inborn preference is accumulated repetition.

\textbf{The second hand: education.}

Schools do not say so outright, but they continually answer the question of what counts as good music. Whom textbooks include or omit, a teacher's tone when mentioning a genre, what choirs sing at competitions: together these form an invisible ranking of taste. In many countries, European classical music---Bach, Mozart---has long occupied the top. These composers are indeed great. But placing them at the top of textbooks is a historical and institutional decision, not the universe's factory setting. Some Indian children receive different training: in Carnatic or Hindustani classical music, a note's intricate turns and a drum pattern's variations can match any symphony in complexity. Yet these traditions rarely appear on other countries' lists of good music.

\textbf{The third hand: money.}

This is the least pleasant hand to discuss, and the hardest to avoid. Piano lessons cost money, live performances require tickets, and a quiet, undisturbed listening environment costs money too. Even having time to explore what you enjoy is, to an extent, something money buys. Compare a child given instrumental lessons and holiday concert visits with one who helps in the family shop after school. How much of their difference in taste at ten is a difference between souls? Asking does not humiliate either child. It brings a hidden ledger of social advantages into the open.

\textbf{The fourth hand: the market.}

You think you select songs, but someone has already filtered the options. Record companies decide whom to promote; platforms decide which songs receive recommendations; short videos may decide a song's success by whether it has a fifteen-second chorus hook. The market watches not your soul but how long you stay. Pause on a song, and it feeds you more like it, drawing a tighter and tighter ring around your preferences as though feeding fish. This is probably how Xiao Ai's viral hit reached her ears.

\textbf{The fifth hand: community.}

People recognize one another through shared likes: matching sneakers, matching profile pictures, knowing the next lyric. Liking something often serves as a community's admission ticket. This is not hypocrisy but a basic human condition: acceptance by a group is a fundamental need. So entire classes develop the same enthusiasms, like an epidemic, or wheat bending together in the wind. Every circle has its must-likes and must-despises; change circles and the list changes completely. Xiao Yu did not stop liking her playlist. Her liking was suddenly exposed as not fitting in.

Having examined all five hands, let us do something crucial: \textbf{make A Zhe's argument as fully as possible}. Looking down on popular taste is often equated immediately with arrogance, which is not entirely fair.

At its strongest, his side has at least three arguments. First, discernment is real. Hearing a difference between songs, like tasting a difference between soups, is an ability developed through time. It deserves respect, just as jumping farther in PE does. Second, serious creators need advocates. If shoddy work always wins, careful creators disappear, and everyone---including Xiao Ai---loses. Third, beware of being spoon-fed. People who never examine their preferences' origins are easiest for businesses and platforms to lead. Some hostility toward the popular can be self-protection in the information age.

This is not a weak case. But you have already spotted its fatal vulnerability: it quietly assumes A Zhe's taste grew independently while Xiao Ai's was spoon-fed. Can that assumption withstand examination? Keep the question with you.

\section{Thinking Bridge: Excavate a Preference}
Feelings alone cannot settle whether your preferences are your own. We need a portable tool. Call it \textbf{taste archaeology}: like an archaeologist excavating a site, dig four layers beneath any ``I like it.''

\textbf{Layer one: exposure---when did I first encounter it?}

Something must enter your world before you can like or dislike it. This layer concerns opportunity: how did this song, food, or kind of film arrive? Who introduced it? Digging often uncovers surprises. You may not remember your first encounter with many supposedly self-chosen preferences, because it happened too early or too routinely. Like the color of a wall that has always been there, it made you assume the world naturally came in that color.

\textbf{Layer two: training---who taught me how to see it?}

Being taught what to notice and receiving no guidance transform the same object into two different experiences. Someone who has watched dozens of games with a coach sees the same match as a first-time spectator but receives different signals. Training need not happen in classrooms: a cousin teaches you to hear the bass line, someone online teaches you to assess sneakers' workmanship, Grandma teaches you to distinguish opera performers. All are training. Ask: who installed the lens through which I appreciate this?

\textbf{Layer three: reward---what does liking it earn in my circle?}

This layer stings most; be honest. Liking a band makes me knowledgeable among certain people. Liking a series lets me join break-time conversations. Hating a kind of song separates me from ``those people.'' This question does not accuse you of hypocrisy: rewards are often unconscious, with no deliberate calculation. But without digging here, you can never distinguish defending the song from defending your place in the group.

\textbf{Layer four: the body---without an audience, do I still like it?}

This is the final shovelful. Imagine being alone, with no phone or social network, and nobody knowing what you hear. Would you still play the song, order the dish, or wear the shoes? If yes, you have uncovered something genuinely rooted in your body: direct pleasure needing no audience. Hesitation does not make you a hypocrite either. It simply means the social share of this preference is larger than you thought.

Try this tool on A Zhe's obscure band. Exposure: his cousin introduced it, the first layer. Training: a forum supplied a whole vocabulary about real music, the second. Reward: a unique classroom position and the attention earned by ``Cleanse your ears,'' the third. The body: he may never have seriously asked himself this fourth question.

After four layers, A Zhe is no more personal than Xiao Ai. That does not mean he loses: her playlist would not withstand these four shovelfuls either. \textbf{Taste archaeology does not pronounce anyone genuine or fake. It does one thing: unpack ``born this way'' into ``how I got here.''} Only when you can see something's origins do you have a chance to become its true master.

\section[Confucius and the Taste of Music]{Confucius's Three Months Without the Taste of Meat}
Now read a short primary source. More than two thousand years ago, someone was overwhelmed by a preference, and left a record of the experience.

The ``Shu Er'' chapter of the Analects records this incident:

\begin{quote}
In Qi the Master heard Shao music, and for three months did not know the taste of meat. He said, ``I never imagined music could reach such heights.''
\end{quote}
In other words, Confucius heard Shao in the state of Qi, an ancient music-and-dance tradition said to praise Shun. For a long time afterward, meat seemed tasteless. He marveled that music could be so beautiful.

Pause here; this passage rewards slow reading. Who was Confucius? One of his age's greatest connoisseurs of music, acquainted with musical theory, the ritual order behind music and dance, and each piece's ceremonial place. In this chapter's terms, his training layer was astonishingly thick. Yet the record describes not analysis but loss of control: three months without the taste of meat. Here was someone deeply trained, utterly submerged in immediate pleasure. Deep understanding and deep enjoyment did not fight in him; they almost seem the same thing. Remember this: we will return to it at the end.

Consider another passage. In ``Gaozi I'' of the Mencius, Mencius uses taste to argue that human nature has common features:

\begin{quote}
Our mouths share preferences in flavors; Yi Ya simply discovered first what our mouths delight in.
\end{quote}
The point is that mouths have shared preferences, and Yi Ya, a famous Spring and Autumn period cook, merely discovered them first. Mencius extends the analogy to ears and eyes: universal admiration for Yi Ya suggests broadly shared tastes, and the same applies to sounds heard by the ears and beautiful sights seen by the eyes.

Put Confucius and Mencius together and you see a confrontation stretching across two thousand years, still unresolved. On Mencius's side, shared human preferences give deliciousness, pleasing sounds, and beautiful sights a common foundation. Ranking taste and judging beauty are therefore not merely domineering; they have grounds. On Confucius's side, at least in this passage, profound aesthetic experience is a private earthquake no ranking can contain. It happens between a particular person and a particular work. Only someone who experienced it could say, ``I never imagined music could reach such heights.''

One argues for commonality; the other demonstrates individuality. Who is right? Do not rush. Together these passages establish at least this: whether taste is personal and whether standards exist are not modern questions invented out of idleness. They are among humanity's oldest questions, and the finest minds have divided over them from the beginning.

\section{One Song, Three Explanations}
Back to the classroom. With more pieces available, we can propose genuinely different explanations for the same facts. First distinguish four things. What happened---Xiao Ai played a song, A Zhe mocked it, Xiao Yu put her phone away---is fact. Why it happened is explanation, where competing accounts clash. How participants understood it---A Zhe felt he defended music, Xiao Ai felt insulted---is their own account. How we evaluate it is a value judgment. We will compare the second layer: \textbf{why this particular song was both liked and despised}.

\textbf{Explanation one: familiarity at work, a psychological account.}

This draws on the psychological pattern already mentioned: repeated exposure produces liking. Xiao Ai likes the song not because it touched her soul, but because the platform presented it dozens of times and her body mistook familiarity for preference. Similarly, A Zhe may not have enjoyed his bands at first. His cousin played them repeatedly until familiarity became appreciation. This explanation is simple, experimentally supported, and evenhanded: both are products of repetition, so neither should despise the other. Its limit is equally clear: it cannot explain differences. Why do you love some repeatedly presented songs and grow increasingly irritated by others? Familiarity is a pass, but it does not open every door.

\textbf{Explanation two: belonging at work, a community account.}

This account looks elsewhere. It asks not whether a song sounds good, but who liking it makes you. Xiao Ai's song is a classroom-wide, internet-wide chorus; liking it buys a ticket to being with everyone. A Zhe likes his band precisely because it is obscure: liking it declares that he is not one of those viral-hit listeners. Here, what is despised is people, not songs; what is defended is boundaries, not music. Xiao Yu withdraws her playlist because its admission ticket has suddenly become more expensive. This sharp explanation captures why disputes over taste are so explosive: they are wars over identity, not ears. But reducing everything to taking sides cannot explain genuinely solitary moments---stumbling across a song alone late at night and feeling your scalp tingle. With no audience present, belonging cannot reach that experience.

\textbf{Explanation three: channels at work, a market account.}

This account looks farther upstream, at supply rather than individuals. Why does Xiao Ai happen to like that song? Because capital selected it: a company funded promotion, a platform provided traffic, short videos required a fifteen-second hook, and producers followed the recipe. Her liking is an assembly line's endpoint. A Zhe has not escaped either: his independent bands are packaged by labels, distributed by platforms, and recommended by algorithms. Being anti-commercial is itself a carefully marketed persona. This explanation follows the money upstream and reminds us that business lies beneath the battlefield of taste. Its limit is a tendency to make people puppets. Yet the assembly line releases ten thousand songs daily and only a few become hits. If taste is entirely spoon-fed, what explains the other nine thousand nine hundred? The market can determine what reaches the table, but not which dish gets finished.

Each explanation holds part of the truth and misses the whole. A familiar story repeats: an account claiming to explain everything has often first discarded something. Familiarity discards differences, belonging discards solitude, and the market discards that final genuine bodily shiver.

\section[Taste as a Classifying Machine]{Deeper Challenge: Taste as a Classifying Machine}

Now for this chapter's weightiest theory, from French sociologist Pierre Bourdieu. In the 1960s and 1970s, he and his colleagues did something then unusual: instead of merely discussing beauty, they entered French homes with questionnaires. They asked about music, paintings, food, furnishings, and leisure, then compared answers with occupations, incomes, and educational qualifications. The results eventually became a substantial book, Distinction, or La Distinction.

Bourdieu's finding can be stated in one sentence: \textbf{taste is distributed across classes, not merely across individuals.} Classical and popular music preferences were not randomly scattered. Broadly, the better educated and more privileged clustered at one aesthetic end, manual workers at another, while those between had their own table manners, holiday practices, and correct preferences. Exceptions existed, but the pattern was as clear as a landscape. Why? Bourdieu offered two key concepts, each worth approaching through everyday examples.

The first is \textbf{cultural capital}. You know what capital does: money can be saved, earn interest, and be passed to children. Culture, Bourdieu argued, can too. A child who hears parents discussing books, visits exhibitions, and has their conversation corrected at dinner accumulates an invisible deposit. They know what to say where, recognize teachers' expectations, and feel unafraid before a painting. In examinations, interviews, and social encounters, that deposit buys real benefits: teachers' instinctive trust, the bearing an interviewer admires, an invitation into a circle. Cultural capital's greatest power is that it \textbf{does not look like money}. Everyone recognizes money as money. Cultural capital appears instead as talent, good breeding, or natural taste. It disguises inherited advantage as individually earned excellence.

The second is \textbf{habitus}. The life you grow up immersed in produces bodily preferences like an accent. Nobody explicitly teaches you which accent you should have; it simply appears when you speak, and you barely recognize it as learned. A worker's child and a doctor's child may carry different bodily memories of a good meal, different feelings about respectable clothing, even different ways of sitting, standing, and laughing. These are products of immersion, not conscious choices. Thus disputes over taste receive a Bourdieusian answer: when A Zhe calls songs assembly-line products, his confidence, vocabulary, and relaxed air of unquestionable authority are themselves products of years in family and school. Xiao Ai's blush enacts another habitus, one that never taught her to defend her taste before the whole class.

Then comes Bourdieu's sharpest incision: \textbf{distinction}. The taste machine's chief product is not pleasure but classification, dividing us from them. ``What I like'' simultaneously means ``whom I am unlike.'' A Zhe's obscure band serves its purpose because others have not heard it. If the whole school listens tomorrow, he must find another---not because the music changed, but because the distinction failed. Hence the endless smoke of taste wars: apparently about beauty, actually defending boundaries.

This theory is powerful, but its qualification must follow immediately to prevent collateral damage. Here is a line this chapter must hold: \textbf{cultural capital does not mean every preference is hypocritical showing off.}

Consider an easily misread counterexample. Jazz emerged in the early twentieth century in small establishments within Black communities in the American South, despised by mainstream society as low-class music. Decades later, it entered concert halls and became a badge of educated middle-class taste; failing to understand it could even mark someone as uncultured. Rock and hip-hop followed almost the same route, from dangerous street-level sounds of the lower classes to music in business students' cars. At first, this seems to refute Bourdieu: lower-class culture can rise, so surely taste has no walls!

Look closer and the opposite appears. Each upward step involved repricing: bars became concert halls, tickets grew expensive, audiences changed, and reviews developed an insider vocabulary. Distinction had not collapsed; it had acquired different merchandise. Nor did those original barroom performers all rise along with jazz's upgraded status. The lesson is not that Bourdieu was wrong, but that the distinction machine can absorb even its opponents.

The opposite misreading also needs guarding against. Some readers conclude that enjoying classical music is mere pretense, that all taste is class performance and nothing is genuine. This pushes the theory too far. Bourdieu explains preferences' \textbf{origins} and \textbf{social functions}; he does not, and cannot, establish that someone's pleasure is false. A working-class child might chance upon opera on the radio, feel overwhelmed, and later scrimp to buy a ticket. That liking still has origins: radio exposure, subsequent training, perhaps a desire to become someone different. Yet the tears shed in the theater are real. Origins cannot invalidate sincerity, any more than knowing a dish's history removes its flavor. \textbf{``Your taste is shaped'' and ``Your pleasure is fake'' are entirely different statements}. The former is almost always right; the latter almost always wrong. Confusing them is an easy trap when reading Bourdieu. It leads people to use theory to humiliate genuine feeling, precisely what the theory most opposes.

\section[Algorithms, Fandom, and Culture Wars]{Algorithms, Fandom, and the ``Junk Culture'' Wars}
Back to today: this ancient war is fighting its latest round in your pocket, with upgraded weapons.

The first new variable is algorithms. Exposure, taste archaeology's first layer, once depended on family, friends, and geography; now much of it belongs to recommendation systems. Their operation fits the familiarity-produces-liking pattern perfectly: linger on a song, and you receive more like it. The more you see something, the more accustomed you become; the more accustomed you become, the more you think you like it. The ring around your preferences tightens faster than ever. This need not be a conspiracy---the system merely wants you to stay longer---but the consequence merits caution. Your personal taste may be a server's by-product, optimized for retention.

The second variable is organized fandom. Admiring an idol can now involve a whole organized way of life: boosting charts, steering comments, buying endorsed products, and defending hashtags. Community's hand has never been stronger. It offers belonging and charges admission; it helps children in distant towns find their people and makes disliking dangerous. Through Bourdieu's eyes, fandom hierarchies---who spends most, has longest-standing credentials, knows more backstage stories---constitute a fresh system of cultural capital, with a new currency.

The third voice is ancient and returns every generation: \textbf{popular culture is doomed.} You hear that disposable pop ruins young people, short videos kill taste, and children can watch only fifteen seconds. Their ancestors said the same. In eighteenth- and nineteenth-century Europe, newspapers and public opinion attacked novels as encouraging escape from reality, moral corruption, and neglect of serious duties, particularly among women and young people. Novels were then today's short videos. Respectable newspapers treated jazz as a menace in the early twentieth century; rock supposedly threatened an entire generation in the 1950s. The script repeats: something new, popular, and beloved by young people appears; those with cultural authority declare it decadent; decades later it becomes a classic used to disparage the next generation's favorite.

That does not make every critique of popular culture wrong. Remember our strongest case for A Zhe. Assembly-line production, calculations about attention time, and shoddy work displacing careful creation are real targets. The key is this: \textbf{criticize the production line, not the ordinary people at its end who genuinely enjoy its products.} Attacking a viral hit's formula is criticism; calling Xiao Ai tasteless for enjoying it is distinction.

Critics often miss popular culture's other half: \textbf{fans have never merely waited open-mouthed to be fed.} Decades ago, Japanese anime and manga fans developed a vast dojin culture of fan-created derivative works. Tens of thousands write, draw, and exchange their creations at conventions large enough to become major city events. Online, fans translate, proofread, and encode subtitles for foreign films and television without pay, simply so more people can watch. Some teach themselves editing, color grading, and music selection to make a three-minute video; others learn a second foreign language to understand an idol's original words. Researchers observing these communities find that fans are often exceptional readers: they analyze frame by frame, investigate details, and offer interpretations more intricate than many professional reviews. In their own terms, they produce as well as consume, recreate as well as receive. Critics describe the public as asleep, but visit an active community and its lights blaze.

Today's picture therefore has two layers. Algorithms and assembly lines are more efficient than ever, making taste's origins more questionable than ever. Yet ordinary people's capacity to participate in and create culture is also stronger than ever, making the word ``audience'' less accurate than ever. Which layer wins depends considerably on which role you, the readers of this book, make a habit of occupying.

\section[Does Appreciation Cost Pleasure?]{For You: Does Learning Appreciation Cost Pleasure?}
Finally, return to those headphones and Confucius's three months.

Imagine a small incident. You have loved a song for years; its introduction always reassures you. One day, a knowledgeable teacher or internet acquaintance analyzes it: the chord progression is thoroughly formulaic, the vocal flourish industrial sweetener, the chorus timed to the millisecond to make you addicted. Is the analysis correct? Quite possibly. But afterward, whenever you listen, a commentator speaks in your ear. The unthinking, scalp-tingling pleasure never returns.

Here, then, is the chapter's question for you, with no standard answer:

\textbf{If learning appreciation might cost you that first immediate pleasure, would you still learn?}

Before answering, count the cards in your hand again.

For learning: discernment is real; those who do not learn are easiest to feed, sell to, and lead. Confucius understood music as deeply as anyone of his time, yet also enjoyed it most profoundly, forgetting meat's taste for three months. Perhaps deep understanding and deep pleasure do not conflict. And when an old pleasure disappears, a greater new one sometimes replaces it: after learning to watch a sport, a match changes from people running around into a battle that makes your palms sweat.

Against learning: some pleasures live only once and cannot be reassembled after dissection. Analysis can become a preparatory course in arrogance. A Zhe knows so much, yet may have no more pleasure than Xiao Ai, only more contempt. Moreover, nobody's understanding of the world is ever complete. If every preference must pass an inspection for sufficient knowledge, people lose the right to simple, wholly unreasoned happiness.

This chapter offers one more card: perhaps the real choice is not whether to learn, but what sort of person to become afterward. Will you use knowledge to extinguish others' pleasure, or to open more windows for yourself without closing theirs?

Give your answer and reasons between Confucius and Mencius, between Xiao Ai and A Zhe, and within the self who hears that commentary. Your playlist remains yours. But from today, you know where its songs came from, and knowing origins is the beginning of every freedom.
